\section{Numerical results}
\label{sec:results}

We simulate a two-level flux-tunable transmon with
$E_C/h=0.2$~GHz, $E_J/h=10$~GHz, $T_1=100~\mu$s, and
$T_2=50~\mu$s.  The reference flux is $\Phi_{\mathrm{opt}}=0.9553\Phi_0$,
the time step is $0.5$~ns, and each $\pi/2$ pulse occupies a 10-ns control
window.  The master equation is solved with QuTiP.  Reported frequency errors
are independently checked against the analytic transmon dispersion, rather
than against the estimate used by the feedback controller.

\subsection{Validity of the local estimator}

Figure~\ref{fig:local-accuracy} compares the linear inverse with two cubic
calibrations while the drive is fixed at the reference frequency.  The full
kernel and odd-polynomial routes agree in the central region, providing an
independent check of the cubic coefficient.  The inferred turnover is
$|\Delta|\simeq17.6$~MHz.  Within the sampled points on this branch, the mean
absolute error is 1.27~MHz for the linear estimator, 0.94~MHz for the fitted
cubic estimator, and 0.97~MHz for the full-kernel estimator.  Near the fold,
the cubic correction can overshoot and the local inverse ceases to improve the
estimate.  The result defines an operating window; it does not justify using
the transient probe for wide-range capture.

\begin{figure*}[t!]
 \includegraphics[width=\textwidth]{transient_accuracy.pdf}
 \caption{Accuracy of the transient frequency discriminator as a function of
 true detuning.  The upper panel shows signed measurement error and the lower
 panel its magnitude.  The shaded interval marks the cubic branch bounded by
 the response turnover.  The fixed-drive polynomial fit and full three-time
 kernel give consistent cubic corrections in the central region.}
 \label{fig:local-accuracy}
\end{figure*}

\subsection{Accuracy--cost tradeoff}

To separate estimator cost from fold failure, we next restrict the flux offsets
to $\pm0.01\Phi_0$, corresponding to $|\Delta|\lesssim12$~MHz.  Figure
~\ref{fig:pareto} compares six strategies.  Double-sweep Ramsey reaches a
6.3-kHz mean absolute error using 800 master-equation solves per frequency
measurement.  Single-sweep Ramsey configurations use 100--400 solves and give
23--30-kHz errors.  The cubic transient estimator gives a 13.5-kHz error at an
amortized cost of 18 solves per point.  Its marginal cost is two solves after a
one-time 80-solve coefficient calibration.  In this fixed-drive benchmark, the
cubic transient point Pareto-dominates all three single-sweep Ramsey settings;
double-sweep Ramsey remains more accurate at substantially higher cost.

\begin{figure*}[t!]
 \includegraphics[width=\textwidth]{accuracy_cost.pdf}
 \caption{Accuracy--cost comparison inside the local transient window.
 (a) Mean absolute frequency error versus measured master-equation calls.
 (b) Error versus wall-clock time on the benchmark machine.  Dashed lines mark
 the nondominated frontier.  Transient cubic cost includes amortization of its
 one-time coefficient calibration over five flux points.}
 \label{fig:pareto}
\end{figure*}

\subsection{Closed-loop calibration}

We set the target frequency 20~MHz below the reference point and use the same
initial flux offset for all strategies.  Figure~\ref{fig:convergence} contrasts
three feedback loops.  A conventional double-sweep Ramsey loop converges in
five iterations.  A coarse transient stage followed by Ramsey refinement takes
seven and five iterations, respectively; the transient stage begins outside
its fold window and wanders before the handoff.  The drive-tracked transient
loop instead converges monotonically in six iterations without a Ramsey
refinement stage.

\begin{figure*}[t!]
 \includegraphics[width=\textwidth]{closed_loop_convergence.pdf}
 \caption{Closed-loop residual for a target shift of $-20$~MHz.  The
 coarse-to-fine workflow (green and blue regions) is compared with the
 drive-tracked full-kernel protocol (purple triangles) and a Ramsey-only
 baseline (red squares).  Residuals are independently evaluated from the
 analytic frequency--flux relation.}
 \label{fig:convergence}
\end{figure*}

Raw iteration count is not a fair cost metric because one Ramsey measurement
contains a delay scan.  We therefore instrument calls to both the master- and
Schrodinger-equation solvers and count each call with unit weight.  This is
conservative for the transient method because the state-vector solver used for
kernel evaluation is cheaper than the master-equation solver.  As shown in
Fig.~\ref{fig:closed-loop-cost}, the drive-tracked loop uses 984 counted calls
and reaches a verified residual of 0.00519~MHz.  Ramsey alone uses 4000 calls
and reaches 0.01797~MHz, giving a 4.1-fold reduction in counted cost.  The
coarse-to-fine workflow uses 5214 calls despite its 0.00214-MHz final residual;
at the common 0.016-MHz tolerance its speedup relative to Ramsey is only 0.8,
so it is more expensive rather than faster.

\begin{figure*}[t!]
 \includegraphics[width=\textwidth]{closed_loop_cost.pdf}
 \caption{Cost anatomy of closed-loop frequency calibration.  (a) Verified
 residual versus cumulative instrumented solver calls.  (b) Cost ratio relative
 to the Ramsey baseline at a common residual threshold.  (c) Total counted
 calls and independently verified final residuals.  The full-kernel transient
 route uses drive tracking; the hybrid route uses a fixed reference in its
 transient stage.}
 \label{fig:closed-loop-cost}
\end{figure*}

The mechanism is visible directly in the measured detuning.  After bootstrap,
drive tracking predicts the next transition frequency closely enough that the
probe detuning falls from sub-megahertz values to a few kilohertz while the
absolute target residual spans tens of megahertz.  The inexpensive probe thus
remains local even though the calibration travels over a nonlocal frequency
interval.  By contrast, keeping the drive fixed forces the probe itself to
span the full target displacement and crosses the turnover identified in
Fig.~\ref{fig:local-accuracy}.
